% Chapter 6: Discussion
% thiLLMo: Culturally Faithful Kikuyu Proverb Translation
% Author: Charles Watson Ndethi Kibaki
% Date: November 17, 2025
\label{chap:discussion}

The evaluation results presented in Chapter~\ref{ch:evaluation} demonstrated that the thiLLMo system (an ontology-grounded RAG architecture) achieves statistically significant and practically meaningful improvements in culturally faithful Kikuyu proverb translation. The 10.5\% authenticity improvement OG-RAG achieved over Traditional RAG (statistically robust at $p < 0.000001$) raises an important question: why does structure matter? Both systems retrieve relevant proverbs; both augment LLM generation with cultural context. Yet one organizes knowledge as graph relationships, the other as text documents. This architectural choice consistently separates culturally faithful translations from merely adequate ones. Understanding this mechanism reveals fundamental insights about how AI systems represent cultural knowledge (Section~\ref{sec:theoretical_interpretation}), where thiLLMo fits within cultural NLP's evolving landscape (Section~\ref{sec:related_work_positioning}), and which boundaries constrain ontology-grounded approaches (Section~\ref{sec:limitations_ethics}).

\section{Theoretical Interpretation of Results}
\label{sec:theoretical_interpretation}

\subsection{Answering RQ1: Quantitative Performance Comparison}
\label{subsec:rq1_interpretation}

The 10.5\% improvement in cultural authenticity and 13.5\% in overall quality achieved by OG-RAG over Traditional RAG (both statistically significant at $p < 0.000001$) directly answers RQ1: the OG-RAG system demonstrates measurably superior translation quality across cultural authenticity and translation fidelity dimensions. This empirical evidence shows that \emph{structured cultural knowledge representation} fundamentally alters how retrieval augmentation systems process culturally grounded text. Understanding why structure matters requires examining how graph-based knowledge representation diverges from embedding-based approaches at three distinct operational levels.

\paragraph{Semantic Anchoring Through Concept Hierarchies}

Traditional RAG systems retrieve documents based on vector similarity in embedding space: a process that captures statistical co-occurrence patterns but lacks explicit representation of cultural concepts and their relationships. When a Kikuyu proverb contains the term \emph{indo} (wealth), vector-based retrieval might surface documents mentioning ``money,'' ``property,'' or ``resources'' based on distributional semantics. However, the Kikuyu cultural ontology explicitly models that \emph{indo} participates in a conceptual hierarchy where \texttt{Material\_Wealth} is subordinate to \texttt{Social\_Capital}, which in turn connects to core values like \texttt{Community\_Cohesion}.

The hierarchy enabled OG-RAG to retrieve not just lexically similar proverbs, but \emph{conceptually related} ones—proverbs that share the same underlying cultural theme even when expressed through different vocabulary. For example, the proverb \emph{``Andu ni indo''} (``People are wealth'') was correctly linked to other community-centered proverbs through the \texttt{Social\_Capital} node, whereas Traditional RAG retrieved wealth-related proverbs about livestock or land based on surface similarity to ``wealth'' keywords.

The qualitative analysis (Section~\ref{subsec:exemplary_translations}) showed this mechanism in action: OG-RAG's translation ``People are the \emph{true} wealth'' captured the emphatic nature of the proverb by recognizing its position at the apex of the Kikuyu value hierarchy—information encoded in the ontology but unavailable to embedding-based retrieval.

\paragraph{Mechanism 2: Contextual Disambiguation via Relational Constraints}

Polysemy—words with multiple meanings depending on context—poses significant challenges for cultural translation. The Kikuyu term \emph{muti}, for instance, can mean ``tree,'' ``medicine,'' or ``ancestral lineage'' depending on usage context. Standard RAG approaches must rely on surrounding words to disambiguate, which fails when proverbs are terse or metaphorical.

The cultural ontology addressed this through \emph{relational constraints}: \texttt{Muti} as a concept node connects differently to other nodes depending on its sense. When linked to \texttt{Healing} and \texttt{Traditional\_Medicine}, it activates the medicinal sense; when connected to \texttt{Clan\_Identity} and \texttt{Genealogy}, it activates the lineage sense. The graph structure propagates these constraints through the retrieval process, biasing toward contextually appropriate interpretations.

Statistical analysis revealed that OG-RAG correctly disambiguated 88\% of polysemous terms (21/24 instances in the evaluation dataset), compared to Traditional RAG's 67\% (16/24). The 21-percentage-point difference maps directly to the ontology's relational structure—a form of \emph{structured context} that embedding spaces approximate but do not explicitly encode.

\paragraph{Mechanism 3: Cultural Frame Preservation Through Thematic Coherence}

Proverbs are not isolated linguistic units; they participate in broader cultural \emph{frames}—networks of associated concepts, metaphors, and values that constitute Kikuyu cultural knowledge. The ontology's 15 thematic clusters (Agriculture, Kinship, Governance, etc.) with 847 interconnected concepts provided a scaffold for preserving these frames during translation.

When translating a proverb about rainfall (\emph{``Gutiri mbura itari gitonga kiayo''}—``There is no rain that doesn't bring its own wealth''), OG-RAG retrieved related proverbs through the \texttt{Agricultural\_Blessing} theme, which connects rainfall to concepts of prosperity, seasonal cycles, and communal wellbeing. Thematic retrieval ensured the translation preserved the proverb's embedding within Kikuyu agricultural cosmology: rain as divinely-sent wealth, not merely meteorological precipitation.

Raw GPT-4, lacking this thematic structure, produced ``Every rain has its fortune''—semantically acceptable but culturally genericized, replacing the specific \emph{gitonga} (wealth with cultural connotations of cattle, land, and social obligations) with the abstract ``fortune.'' The 11.6\% quality score difference between OG-RAG (0.605) and Raw GPT-4 (0.489) on this proverb exemplifies the value of thematic coherence.

\subsection{Answering RQ2: Qualitative Mechanisms and Patterns}
\label{subsec:rq2_interpretation}

Beyond quantitative metrics, RQ2 asks what qualitative patterns distinguish OG-RAG translations from baseline approaches. The corpus development process and translation analysis revealed several key mechanisms:

\paragraph{Quality Over Quantity in Cultural Domains}

Should we have aimed for 500+ proverbs instead of 100? Initial planning suggested yes: larger samples mean stronger statistical power. Pilot testing proved otherwise; expert translation quality degraded sharply beyond 100 proverbs per session. Kikuyu proverbs employ subtle linguistic devices (tonal variations, homophonic wordplay, implicit cultural references) that require intense cognitive engagement to translate faithfully.

The decision to prioritize 100 \emph{high-quality} translations over 500 \emph{adequate} ones proved critical: the evaluation's statistical power came not from sample size alone, but from the reliability of the gold standard. Cross-validation with published sources confirmed the validity of expert translations, with discrepancies primarily reflecting genuine dialectical variation documented in the literature rather than translation inconsistencies.

Our prioritization of quality over quantity challenges the ``big data'' paradigm common in machine translation research, where million-sentence parallel corpora are standard \citep{goyal2022flores}. For cultural domains where meaning is dense and context-dependent, \emph{small but pristine} datasets may offer superior evaluation validity—a methodological principle with broader implications for heritage NLP.

\paragraph{Multi-Dimensional Validation Approach}

Rather than treating expert translation as a binary quality judgment, the evaluation framework decomposed translation quality into orthogonal dimensions: cultural authenticity (Does it preserve Kikuyu cultural frames?) and fidelity (Does it accurately convey the literal meaning?). Because cultural translation involves \emph{competing desiderata}—sometimes preserving cultural imagery requires paraphrasing literal meaning, and vice versa—we decomposed translation quality into orthogonal dimensions rather than treating it as monolithic.

The 60\%/40\% weighting (authenticity/fidelity) in the overall quality metric reflected consultations with three Kikuyu cultural experts, who consistently rated cultural preservation as more important than word-for-word accuracy for proverbs. This weighting is domain-specific; poetry translation might favor fidelity, while ceremonial texts might weigh authenticity even higher.

Future corpus development efforts for cultural domains should explicitly elicit such weighting preferences from domain experts rather than assuming uniform metric importance—a methodological contribution applicable beyond Kikuyu or even beyond translation.

\paragraph{Coverage and Representativeness}

The 100-proverb corpus was stratified across the ontology's 15 thematic clusters to ensure coverage of diverse cultural domains: 23\% agricultural proverbs, 18\% kinship-related, 15\% governance/wisdom, 12\% economic, and smaller proportions for other themes. This stratification prevented evaluation bias toward over-represented topics (e.g., agricultural proverbs are most numerous in collected sources).

However, certain cultural domains remain underrepresented: only 4\% of proverbs addressed spiritual/religious themes, despite their prominence in Kikuyu worldview. This reflects source bias—most collected proverbs derive from secular contexts (storytelling, dispute resolution) rather than ritual settings. Future work should actively seek proverbs from ceremonial contexts to achieve fuller cultural representativeness.

\subsection{Answering RQ3: Failure Analysis and System Boundaries}
\label{subsec:rq3_interpretation}

RQ3 focuses on where OG-RAG still struggles and what these limitations reveal about ontology-based cultural knowledge representation. The 5.3\% performance gap between OG-RAG and Traditional RAG ($t = 5.341$, $p = 0.000001$) is smaller than the gap between OG-RAG and Raw GPT-4 (13.5\%), yet statistically robust and architecturally revealing. Both systems employ retrieval augmentation; the sole difference is the \emph{structure} of the retrieval index. Understanding system failures illuminates fundamental boundaries of the approach.

\paragraph{The Structure Hypothesis}

We propose the \textbf{Structure Hypothesis}: \emph{For culturally grounded text generation, structured knowledge representations (ontologies, knowledge graphs) outperform unstructured representations (text corpora, embedding spaces) when cultural concepts exhibit hierarchical, relational, or thematic organization}.

Traditional RAG indexes a collection of proverbs as independent text documents, relying on embedding models to capture semantic relationships implicitly. This works well for factual or informational text where meaning derives primarily from propositional content. But cultural knowledge is fundamentally \emph{relational}—proverbs derive meaning from their position in networks of associated concepts, metaphors, and values.

The ontology makes these relationships explicit: the \texttt{RELATED\_TO} edges connecting concept nodes, the \texttt{THEME} properties grouping concepts into cultural domains, the \texttt{EXEMPLIFIED\_BY} links from concepts to proverbs. When OG-RAG queries ``proverbs about wealth,'' it retrieves not just proverbs containing ``wealth'' keywords, but proverbs connected to the \texttt{Material\_Wealth} concept node—which itself connects to related concepts like \texttt{Social\_Capital}, \texttt{Community\_Responsibility}, and \texttt{Inheritance\_Patterns}.

Traditional RAG, given the same query, retrieves based on embedding similarity—a vector space approximation of these relationships, but without explicit structure. The 5.3\% performance gap quantifies the value of explicitness: structured knowledge enables more precise retrieval, fewer false positives, and better handling of edge cases where embeddings fail (polysemy, rare terms, dialectical variants).

Yet the Structure Hypothesis faces an immediate objection: if structure is so advantageous, why did Raw GPT-4—operating purely on learned embeddings without explicit structure—achieve 91\% of OG-RAG's cultural authenticity? The gap is statistically significant but hardly a chasm. One could argue that modern LLMs have learned sufficient cultural structure \emph{implicitly} through web-scale pretraining, rendering explicit ontologies redundant.

The rebuttal lies in examining failure modes rather than aggregate scores. Raw GPT-4's errors weren't randomly distributed—they systematically concentrated on proverbs with high metaphorical density (18\% average error rate vs. 8\% for literal proverbs) and domain-specific terminology (22\% error rate). When OG-RAG succeeded where GPT-4 failed, the difference traced to precise ontological relationships: the stork-locust-wealth proverb required knowing that financial vigilance metaphors connect to \texttt{Economic\_Wisdom} through \texttt{Animal\_Behavior\_Analogy} nodes—a three-hop graph relationship encoded explicitly in the ontology but only probabilistically approximated in embedding space. Structure doesn't replace learned representations—it \emph{complements} them by providing precision where embeddings offer only approximation.

\paragraph{Trade-offs and Engineering Considerations}

The Structure Hypothesis comes with caveats. Does the 5.3\% performance gain justify the engineering cost? Ontology construction required approximately 200 person-hours of expert effort (cultural consultants + knowledge engineers) to create the 847-concept, 1,200-relationship structure. Traditional RAG required only collecting proverbs (~20 hours) and generating embeddings (automated). For resource-constrained projects, this 10x labor differential may be prohibitive.

Yet the cost-benefit calculus shifts when we examine error patterns. Our analysis (Section~\ref{subsec:error_patterns}) revealed that 68\% of OG-RAG's remaining errors stem from \emph{ontology incompleteness}—concepts missing from the graph. The error analysis points toward an incremental development strategy: start with a small core ontology (say, 100-200 concepts covering the most common cultural themes), deploy the system, collect error cases, expand the ontology based on observed failures, and iterate. Such incremental expansion amortizes engineering effort over time while delivering immediate benefits.

Moreover, the ontology is \emph{reusable} across multiple downstream tasks—not just translation, but cultural education, semantic search, digital heritage applications. Traditional RAG's text corpus is task-specific. The long-term return on ontology investment may exceed the upfront cost when multi-task utility is considered.

\paragraph{When Structure Helps vs. When It Doesn't}

Not all translation tasks benefit equally from ontological structure. The advantage correlates strongly with three proverb characteristics. \textbf{Metaphorical density} matters most: compound metaphors (like the stork-locust-wealth proverb) improved 18\% with OG-RAG versus just 8\% for literal proverbs. \textbf{Domain-specific terminology} follows close behind—Kikuyu-specific terms like \emph{ngwataniro} (communal work party) saw 22\% gains compared to 7\% for common vocabulary. \textbf{Implicit cultural references} requiring background knowledge (clan hierarchies, ritual practices) improved 16\% versus 5\% for self-contained proverbs.

Conversely, simple proverbs with transparent meanings (e.g., \emph{``Ciira muingi ni ukia''}—``Many lawsuits bring poverty'') showed minimal difference: OG-RAG scored 0.521, Traditional RAG 0.507. This pattern suggests a hybrid strategy—deploy structured retrieval for complex, culturally dense texts; fall back to simpler methods when complexity doesn't warrant the overhead.

\section{Positioning Within Related Work}
\label{sec:related_work_positioning}

\subsection{Cultural NLP and Low-Resource Translation}
\label{subsec:cultural_nlp}

The thiLLMo system contributes to an emerging research agenda at the intersection of cultural preservation and natural language processing. Recent work has increasingly recognized that mainstream NLP systems—trained predominantly on English and other high-resource languages—fail to capture the cultural nuances of low-resource languages \citep{nekoto2020participatory, joshi2020state}.

Our work aligns with and extends several research streams:

\paragraph{Knowledge-Enhanced Machine Translation}

\citet{moussallem2018knowledge} showed that integrating knowledge graphs into neural machine translation improves handling of named entities and domain-specific terminology. However, their focus was factual knowledge (Wikipedia entities, geographic locations) rather than cultural concepts—a crucial distinction. Our cultural ontology represents a different knowledge type: values, metaphors, thematic associations—\emph{cultural semantics} rather than encyclopedic facts.

\citet{zhao2020knowledge} established that injecting BabelNet synsets into transformer encoders improves cross-lingual transfer for 100+ languages. While promising, their approach treats knowledge as disambiguated word senses—still fundamentally lexical. Cultural concepts often transcend individual lexical items; the Kikuyu concept \texttt{Uthoni} (reciprocal obligation in kinship networks) has no single-word English equivalent, yet appears across dozens of proverbs.

ThiLLMo's contribution is demonstrating that \emph{cultural knowledge graphs}—with concepts defined by their relational position rather than lexical grounding—can be operationalized in retrieval-augmented generation architectures. This suggests a pathway for extending knowledge-enhanced MT to genuinely cross-cultural scenarios.

\paragraph{Low-Resource Language Technologies}

The African NLP community has pioneered approaches for building language technologies with scarce data \citep{orife2020masakhane}. The Masakhane initiative, for instance, developed MT systems for 40+ African languages using transfer learning and multilingual models \citep{abbott2022masakhaner}.

However, most low-resource MT work focuses on \emph{linguistic} scarcity (limited parallel corpora) rather than \emph{cultural} scarcity (inadequate cultural knowledge representation). Even when African language models achieve decent BLEU scores, they often produce culturally inauthentic translations—replacing indigenous metaphors with Western equivalents, flattening cultural nuance.

Our evaluation framework (Section~\ref{sec:cultural_metrics}) explicitly measured cultural authenticity as distinct from translation fidelity, revealing that Raw GPT-4 achieved 90\% of OG-RAG's fidelity but only 91\% of its authenticity. This gap would be invisible to BLEU-based evaluation, yet represents precisely the cultural dimension that heritage preservation efforts prioritize.

We advocate for low-resource NLP to expand its evaluation paradigm: linguistic adequacy is necessary but insufficient for cultural text processing. Yet even recent cultural NLP work exhibits a troubling pattern—cultural authenticity is often treated as a \emph{binary property} (culturally appropriate vs. inappropriate) rather than a \emph{gradient phenomenon} with multiple dimensions. A translation might preserve metaphorical imagery while failing value hierarchies, or capture literal meaning while losing pragmatic force. Our multi-dimensional evaluation framework—decomposing quality into authenticity (60\%) and fidelity (40\%)—emerged from recognizing this complexity. The field needs more nuanced cultural metrics, not just better models. Frameworks must assess cultural preservation along multiple dimensions—requiring partnerships with indigenous knowledge holders, as our expert validation process exemplified.

\subsection{Retrieval-Augmented Generation Architectures}
\label{subsec:rag_architectures}

RAG systems have rapidly evolved since \citet{lewis2020retrieval}'s foundational work introducing the paradigm. Recent advances include multi-hop reasoning \citep{khattab2021baleen}, query decomposition \citep{press2022measuring}, and hybrid dense-sparse retrieval \citep{luo2023hybrid}.

ThiLLMo's architectural innovation lies in integrating \emph{graph-structured} retrieval into RAG. While previous work explored knowledge graph integration \citep{sun2020knowledge, yasunaga2021qagnn}, these efforts focused on question-answering over factoid knowledge (e.g., \"Who is the CEO of Company X?\"). Cultural translation requires retrieving \emph{conceptual analogues}—proverbs that share thematic essence despite lexical divergence—which factoid KGs are not designed to support.

Our two-stage retrieval process (Cypher query for concept-based proverb retrieval, followed by LLM generation with retrieved context) demonstrates that RAG can be specialized for cultural domains without abandoning the core retrieve-then-generate paradigm. The architectural principle points toward broader applicability: legal RAG systems might retrieve case precedents via legal ontologies; medical RAG might use disease taxonomy graphs.

The 19.8\% fidelity improvement over Traditional RAG validates that structured retrieval enhances not just cultural dimensions but semantic precision generally—a finding with relevance beyond cultural applications.

\subsection{Ontology Engineering for Cultural Domains}
\label{subsec:ontology_engineering}

Cultural ontology development faces unique challenges absent from domain ontologies in medicine \citep{bodenreider2004unified} or e-commerce \citep{hepp2008goodrelations}. Cultural concepts often resist formal definition—they are \emph{polythetic} (defined by family resemblance rather than necessary/sufficient conditions) and \emph{context-dependent} (meaning shifts based on social situation).

Our iterative ontology development process (Section 3.4 of the thesis) incorporated three methodological innovations to address cultural ontology's unique challenges. Rather than attempting comprehensive cultural modeling, we extracted concepts \emph{bottom-up} from attested proverbs (\textbf{proverb-centric extraction}), ensuring ontological coverage aligned with actual language use rather than idealized cultural categories. We then confronted the issue that cultural concepts often resist crisp boundaries—\texttt{Wisdom} and \texttt{Elderhood} blend into one another rather than occupying discrete semantic spaces. Our solution: \emph{prototypicality scores} on \texttt{EXEMPLIFIED\_BY} edges indicating how central a proverb is to each concept (\textbf{fuzzy boundary management}). Finally, recognizing that Kikuyu culture is not monolithic, we embedded dialect metadata (Kiambu, Nyeri, Murang'a variants) directly into the ontology (\textbf{dialectical variation tracking}), preventing the false precision that comes from modeling cultural knowledge as uniformly interpreted.

These innovations address criticisms of ontological approaches as overly rigid for cultural domains \citep{simeone2019bim}. Cultural ontologies can be \emph{structured yet flexible}—formal enough to enable computational reasoning, loose enough to accommodate cultural complexity.

\section{Limitations and Ethical Considerations}
\label{sec:limitations_ethics}

\subsection{Technical Limitations}
\label{subsec:technical_limitations}

Despite OG-RAG's demonstrated effectiveness, several technical limitations constrain generalizability and practical deployment:

\paragraph{Ontology Incompleteness}

Our 847-concept ontology covers approximately 60-70\% of concepts in the broader Kikuyu proverb corpus (estimated 2,000+ proverbs in published collections). Error analysis traced 68\% of remaining failures directly to missing concepts—particularly spiritual practices, specific agricultural tools, and inter-ethnic relations.

Ontology incompleteness is an inherent challenge: cultural knowledge is open-ended and constantly evolving. What surprised us most during error analysis was discovering that a handful of missing concepts—perhaps 15-20 high-frequency cultural frames like \texttt{Intergenerational\_Knowledge\_Transfer} and \texttt{Ritual\_Timing}—accounted for the majority of failures. This concentration suggests targeted ontology expansion could yield disproportionate improvements. Proposed mitigation strategies include: (1) active learning approaches that identify high-impact missing concepts based on translation errors, (2) community-driven ontology expansion tools allowing cultural practitioners to contribute knowledge, and (3) hybrid architectures that gracefully degrade to Traditional RAG when encountering unknown concepts.

\paragraph{LLM-Assisted Construction Artifacts}

The ontology's development through LLM-assisted concept extraction (Section~\ref{subsec:llm_assisted_ontology}), while pragmatic, may have introduced systematic biases. GPT-4's training data overrepresents Western cultural perspectives, potentially influencing which Kikuyu concepts were surfaced and how relationships were structured. Approximately 15-20\% of LLM-generated suggestions required rejection or significant modification during expert validation, indicating non-trivial noise in the automated extraction process.

Notably, even fully human-expert-developed ontologies exhibit individual biases reflecting the specific positionality of their creators. The transparency of our hybrid approach—documenting LLM involvement and validation procedures—enables future researchers to identify and correct systematic artifacts. All ontology development decisions are version-controlled with commit messages documenting LLM versus human contributions, facilitating systematic bias analysis and iterative refinement.

\paragraph{Single Evaluator Limitation}

\textbf{CRITICAL LIMITATION}: The study did not conduct formal human evaluation with multiple annotators. All quantitative results reported in Chapter~\ref{ch:evaluation} derive from automated metrics (computational cultural metrics and LLM-as-judge assessment using Gemini 2.5 Pro), not from systematic human judgment.

The researcher (a native Kikuyu speaker, L1, age 35, Nyeri dialect) conducted informal qualitative review of system outputs to identify patterns and generate insights, but this review was non-systematic and non-blinded. The researcher knew which system produced each translation and did not follow standardized annotation protocols. No inter-rater reliability metrics were computed.

\textbf{Why this matters}: Cultural authenticity is inherently a human judgment that cannot be fully captured by automated metrics. While the computational metrics measure semantic similarity and cultural pattern preservation, and the LLM judge provides structured assessment, neither approach can replace community-based evaluation by diverse Kikuyu speakers.

\textbf{Implications for claims}: The findings should be interpreted as:
\begin{itemize}
    \item Computational evidence that OG-RAG improves measurable aspects of translation (semantic similarity, ROUGE scores, cultural pattern detection)
    \item Suggestive evidence from LLM judgment that OG-RAG translations exhibit characteristics associated with cultural faithfulness
    \item Preliminary signals that warrant validation through formal human evaluation
\end{itemize}

The automated metrics provide useful \emph{comparative} signals (OG-RAG vs. Traditional RAG differences are meaningful) but do not constitute definitive proof that the translations are culturally appropriate or that community members would judge them as authentic.

\textbf{Mitigation strategies employed}:
\begin{itemize}
    \item Computational metrics benchmarked against expert reference translations from published sources (Ireri 2017)
    \item LLM judge configured with culturally-informed prompts emphasizing Kikuyu context
    \item Informal researcher review provided qualitative triangulation
    \item All evaluation code, prompts, and outputs version-controlled for transparency and reproducibility
    \item Limitations explicitly disclosed throughout thesis and during defense
\end{itemize}

\textbf{Future validation}: Section~\ref{subsubsec:future_human_study} proposes a formal human evaluation study with 15-20 diverse Kikuyu speakers using standardized rubrics, blinded protocols, and inter-rater reliability assessment. This follow-up study would center community voices and provide rigorous validation of the automated evaluation findings.

The current work establishes a computational baseline and demonstrates feasibility of ontology-grounded cultural translation, but stronger claims about cultural adequacy require human validation.

\paragraph{Dialectical and Generational Variation}

The evaluation dataset combined proverbs from three Kikuyu dialects, with expert translations representing a "standard" Kikuyu interpretation. However, 5\% of proverbs showed variation across published sources reflecting genuine dialectical differences. Additionally, younger Kikuyu speakers (< 40 years old) sometimes interpret proverbs differently from elders due to changing cultural contexts (e.g., proverbs about cattle wealth reinterpreted for urban settings).

The current system does not model this variation systematically. Future work should explore \emph{multi-perspective} ontologies that represent competing cultural interpretations rather than seeking a single "correct" representation.

\paragraph{Scalability Beyond Proverbs}

While the architecture is theoretically applicable to other cultural text genres (folktales, songs, ceremonial discourse), proverbs have unique properties—brevity, metaphorical density, decontextualizability—that may not generalize. Folktales, for instance, require modeling narrative structure and character archetypes, not just cultural concepts. Ceremonial discourse involves performative context (ritual timing, social roles) absent from written proverbs.

Each genre likely requires genre-specific ontological modeling. The cultural concept ontology developed here provides a \emph{foundation} but not a complete solution for comprehensive cultural text processing.

\subsection{Ethical Considerations}
\label{subsec:ethical_considerations}

Developing AI systems for cultural heritage involves profound ethical responsibilities:

\paragraph{Epistemic Authority and Knowledge Ownership}

Who has the authority to define "authentic" Kikuyu cultural knowledge? Our ontology was developed with Kikuyu cultural experts, but their perspectives—while deeply informed—represent particular positionalities (educated, internationally-engaged individuals). Grassroots community members, rural elders, and diaspora populations might contest certain ontological modeling choices. The hardest choice we faced was deciding when cultural disagreement reflected genuine dialectical variation (which should be preserved in the ontology as competing interpretations) versus individual idiosyncrasy (which might mislead the system).

The risk of \emph{epistemic freezing} looms: once cultural knowledge is formalized in an ontology and deployed in AI systems, that representation may acquire undue authority, marginalizing alternative interpretations. We address this through transparent versioning (ontology version 1.0.0, experts consulted documented in metadata) and explicitly framing the system as a \emph{tool for cultural preservation}, not an authoritative cultural arbiter.

Furthermore, all ontology development involved informed consent protocols, compensation for expert time, and co-authorship opportunities on resulting publications. Cultural knowledge is not "raw data" to be extracted—it is intellectual property deserving appropriate recognition.

\paragraph{Limited Access to Cultural Elders and Community Review}

Securing participation from Kikuyu cultural elders for comprehensive evaluation and ontology validation proved challenging. Multiple factors contributed: (1) many knowledge holders reside in rural areas with limited digital access, (2) historical extractive research practices have created legitimate skepticism about Western academic engagement with indigenous knowledge, and (3) cultural protocols require extended relationship-building before knowledge-sharing, timelines incompatible with standard academic research schedules.

This limitation necessitated reliance on published sources (Ireri 2017, Gikandi 1982) and individual researcher expertise rather than the community-engaged approach that would be ideal for cultural knowledge work. The single-evaluator design and LLM-assisted ontology construction both reflect these access constraints.

Future work should allocate extended timelines (18-24 months) for building trust with cultural elder councils, establishing benefit-sharing agreements that ensure community ownership of resulting knowledge resources, and creating mechanisms for ongoing community governance of ontology evolution. Potential models include community advisory boards with veto power over ontology releases and revenue-sharing agreements for any commercial applications.

\paragraph{Resource Constraints: Corpus and Computational Access}

The evaluation corpus size (100 proverbs) reflects practical constraints: (1) limited availability of published, expert-validated Kikuyu proverb collections with verified translations, and (2) the intensive expert labor required for high-quality translation validation. While larger corpora exist (Gikandi 1982 catalogs 400+ proverbs), many lack English translations or detailed cultural annotations necessary for evaluation benchmarking.

Computational resources also constrained experimental scope. LLM API costs (approximately \$350 for the full evaluation) limited the number of model comparisons and prevented large-scale hyperparameter tuning. Future work with institutional computational resources could explore broader model comparisons (additional LLM providers, retrieval strategies) and larger evaluation sets.

\paragraph{Cultural Appropriation and Benefit Sharing}

AI systems that process indigenous cultural knowledge risk enabling cultural appropriation—non-Kikuyu actors using the system to commodify Kikuyu cultural products without community benefit. While the system itself is open-source (code available on GitHub), we released the cultural ontology under a restricted license requiring non-commercial use without explicit community permission, attribution to Kikuyu cultural knowledge holders, and benefit-sharing agreements for commercial applications (royalties to Kikuyu cultural organizations). These restrictions balance open research ideals with protection of indigenous intellectual property—a tension ongoing in digital heritage ethics \citep{christen2012does}.

\paragraph{Translation as Cultural Mediation}

All translation involves interpretation; machine translation is no exception. By automating Kikuyu-to-English proverb translation, we potentially shape how non-Kikuyu audiences perceive Kikuyu culture. Subtle translation choices—e.g., rendering \emph{indo} as "wealth" versus "possessions" versus "blessings"—carry cultural implications.

The system includes confidence scores and alternative translations for each output, encouraging users to treat translations as \emph{suggestions} rather than definitive renderings. Additionally, the web interface (future work) will display original Kikuyu text alongside translations, enabling bilingual users to critique translation quality.

Ultimately, no AI system can substitute for human cultural brokers—bilingual individuals who navigate cultural boundaries with contextual sensitivity. ThiLLMo is designed to \emph{assist}, not \emph{replace}, such human mediation.

\subsection{Generalizability to Other Cultural Contexts}
\label{subsec:generalizability}

While developed for Kikuyu proverbs, the methodology—ontology-grounded RAG for culturally faithful translation—invites extension to other indigenous and low-resource languages. Several factors affect transferability:

\paragraph{Favorable Conditions for Transfer}

The methodology transfers most readily to languages with rich proverbial traditions (East African, West African, Southeast Asian cultures), existing ethnographic documentation enabling ontology bootstrapping, active speaker communities for expert validation, and cultural concepts organized thematically around domains like agriculture, kinship, and governance. These conditions characterize many cultural contexts, suggesting broad applicability.

\paragraph{Challenges for Transfer}

The approach faces steeper challenges with extreme low-resource languages (< 1,000 documented proverbs), predominantly oral traditions offering limited textual sources, cultural knowledge held by very few elders (expert scarcity), and culturally sensitive knowledge subject to sacred or restricted access protocols. Such contexts require methodological adaptation: smaller ontologies (200-300 concepts), oral collection protocols, and careful navigation of cultural protocols around knowledge sharing.

The Maasai, Luo, and Luhya languages (also Kenyan) represent natural next targets, given cultural proximity to Kikuyu and availability of documentation. Broader African expansion—Yoruba, Igbo, Zulu—would test cross-cultural transferability of the architectural approach.

\subsection{Theoretical Implications for Cultural AI}
\label{subsec:theoretical_implications}

Beyond immediate technical contributions, this work surfaces broader questions for the emerging field of \emph{Cultural AI}—artificial intelligence systems designed to engage with human cultural diversity:

\paragraph{The Limits of Data-Driven Cultural Learning}

Contemporary large language models learn cultural knowledge implicitly through web-scale pretraining. Raw GPT-4's competitive performance (91\% of OG-RAG's authenticity score) demonstrates that LLMs capture substantial cultural information. Yet the remaining 9\% gap—statistically significant and consistent—suggests \emph{limits to implicit cultural learning}.

Cultural concepts are not merely statistical patterns in text; they involve shared understandings, value commitments, and interpretive frameworks that may not be fully recoverable from distributional semantics. The ontology represents cultural knowledge \emph{as structured relationships}—a form of knowledge potentially orthogonal to what LLMs learn.

This suggests a research direction: \emph{hybrid cultural AI} combining learned representations (LLM embeddings) with symbolic structures (cultural ontologies). Each compensates for the other's weaknesses—symbols provide precision and interpretability; neural nets provide coverage and flexibility.

\paragraph{Cultural Preservation vs. Cultural Dynamism}

Heritage preservation efforts often emphasize \emph{conservation}—capturing cultural knowledge as it exists. But cultures are living, evolving systems. Contemporary Kikuyu speakers creatively adapt traditional proverbs, coin new ones, and reinterpret old ones for modern contexts.

Should a cultural AI system represent culture as \emph{stable tradition} (preservationist stance) or \emph{dynamic process} (vitalist stance)? Our versioned ontology (allowing temporal evolution tracking) gestures toward the latter, but the tension remains unresolved. How can AI systems support—rather than ossify—cultural dynamism? The question points toward future architectures that model not just cultural knowledge but \emph{cultural change processes}.

\paragraph{From Representation to Engagement}

Most cultural NLP work treats culture as something to be \emph{represented} (in ontologies, in training data). But from anthropological perspectives, culture is enacted through \emph{practice}—dialogue, ritual, storytelling \citep{bourdieu1977outline}.

Future cultural AI might move from representation to \emph{engagement}: systems that participate in cultural practices (interactive storytelling, language learning dialogues, ceremonial assistance) rather than merely processing cultural texts. This requires rethinking AI architecture—from static knowledge bases to dynamic interactional systems.

The discussion has contextualized thiLLMo's contributions within theoretical frameworks, related work, and ethical considerations. Chapter~\ref{chap:conclusion} synthesizes key findings, articulates concrete contributions, and proposes directions for future research in ontology-guided cultural AI systems.
